% !TEX TS-program = luatex
% Awesome Source CV LaTeX Template
%
% This template has been downloaded from:
% https://github.com/darwiin/awesome-neue-latex-cv
%
% Author:
% Christophe Roger (Darwiin)
%
% Template license:
% CC BY-SA 4.0 (https://creativecommons.org/licenses/by-sa/4.0/)

\documentclass[localFont,alternative,10pt,compact]{yaac-another-awesome-cv}
\name{Samuel}{SIKATI KENMOGNE}
\tagline{Ingénieur Logiciel, Assurance Qualité \& Analyse Fonctionnelle}
\socialinfo{
	\smartphone{(+237) 657 407 989}
	\email{sikatikenmogne@gmail.com}
	\linkedin{samuelsikati}
	\github{sikatikenmogne}
	\website{https://samuelsikati.com}{samuelsikati.com}
}
\headeraddress{Douala, Cameroun}

% Métadonnées PDF (lues par les ATS et les moteurs de recherche)
\hypersetup{
	pdftitle    = {CV - Samuel SIKATI KENMOGNE - Ingénieur QA / Test Logiciel},
	pdfauthor   = {Samuel SIKATI KENMOGNE},
	pdfsubject  = {Curriculum Vitae - Ingénieur QA, Test Logiciel et Assurance Qualité},
	pdfkeywords = {Ingénieur QA, Ingénieur Test, Testeur, Test logiciel, Assurance Qualité Logicielle, Recette fonctionnelle, Homologation, UAT, Cahier de recette, Stratégie de test, Tests fonctionnels, Tests d'intégration, Tests E2E, Tests API, Tests de régression, Tests de performance, Migration de données, PRA, Automatisation des tests, Playwright, Locust, ISTQB, Bruno, Postman, pytest, jest, Squash TM, BDD, Gherkin, Gestion des anomalies, SQL, CI/CD, Jenkins, GitLab, Azure DevOps, SonarQube, Mobile Money, USSD, Scrum, Douala, Cameroun},
	pdfcreator  = {LaTeX},
	pdflang     = {fr-FR},
	pdfdisplaydoctitle = true
}

% Uncomment the following line and use a value from 1.5cm to 2.5cm
% \setleftcolumnlength{2.5cm}

% Décommenter pour ajouter la bibliographie (voir section_publications.tex)
%\addbibresource{my_publications.bib}

\begin{document}

	\makecvheader

	\makecvfooter[{\itshape Fait à Douala, le \today. Je certifie sur l'honneur l'exactitude des informations contenues dans ce CV.}]
		{\textsc{\today}}
		{\textsc{CV - Ingénieur QA / Test Logiciel}}
		{Page \thepage / \pageref*{LastPage}}

	% Awesome Source CV LaTeX Template
%
% This template has been downloaded from:
% https://github.com/darwiin/awesome-neue-latex-cv
%
% Author:
% Christophe Roger
% Contenu adapté pour : Samuel SIKATI KENMOGNE
%
% Template license:
% CC BY-SA 4.0 (https://creativecommons.org/licenses/by-sa/4.0/)

%Section: Accroche
\par{
	Jeune ingénieur logiciel avec une solide expérience pratique en développement d'applications métiers, passionné par l'innovation technologique et la transformation digitale.
}
				% Accroche
	% Awesome Source CV LaTeX Template
%
% This template has been downloaded from:
% https://github.com/darwiin/awesome-neue-latex-cv
%
% Author:
% Christophe Roger
% Contenu adapté pour : Samuel SIKATI KENMOGNE
%
% Template license:
% CC BY-SA 4.0 (https://creativecommons.org/licenses/by-sa/4.0/)

%Section compétences
\sectionTitle{Compétences}{\faTasks}

	\begin{keywords}
		\keywordsentry{Assurance Qualité logicielle}{Recette fonctionnelle (UAT), tests E2E et de régression, BDD (\textbf{Gherkin}), gestion des tests (\textbf{Squash TM}), tests API (\textbf{Bruno}, \textbf{Postman}), automatisation (\textbf{Playwright}, pytest, jest)}
		\keywordsentry{Analyse fonctionnelle}{Analyse des exigences, processus métier (\textbf{BPMN}), modélisation des systèmes (\textbf{UML}, \textbf{SysML})}
		\keywordsentry{Bases de données}{\textbf{SQL}, \textbf{PostgreSQL}, \textbf{SQL Server}, MongoDB, Supabase}
		\keywordsentry{Développement Full Stack}{\textbf{.NET Core} (C\#), \textbf{Symfony} (PHP) \quad | \quad API REST (JWT), SOAP \quad | \quad \textbf{React}, Tailwind CSS}
		\keywordsentry{DevOps \& Méthodes}{Git, GitLab / Azure DevOps, Jenkins, SonarQube, \textbf{Docker}, Linux \quad | \quad Scrum, Kanban}
	\end{keywords}
				% Section compétences
	% Awesome Source CV LaTeX Template
%
% This template has been downloaded from:
% https://github.com/darwiin/awesome-neue-latex-cv
%
% Author:
% Christophe Roger
% Contenu adapté pour : Samuel SIKATI KENMOGNE
%
% Template license:
% CC BY-SA 4.0 (https://creativecommons.org/licenses/by-sa/4.0/)

%Section: Work Experience at the top
\sectionTitle{Expériences Professionnelles}{\faSuitcase}
%\renewcommand{\labelitemi}{$\bullet$}
\begin{experiences}
  \experience
    {Nov. 2026}   {Ingénieur SI / Analyste en Assurance Qualité Logicielle}{Orange Cameroun}{Douala}
    {Nov. 2025} {
                      \begin{itemize}
                        \item Conception, exécution et automatisation des tests fonctionnels et d'acceptation (UAT) dans le cadre de la \textbf{migration d'une plateforme Mobile Money}
                        \item \textbf{Validation end-to-end des flux transactionnels critiques}, exécution des campagnes de tests et qualification des anomalies
                        \item Formation des utilisateurs métiers
                        \item Configuration, paramétrage et recette fonctionnelle d'une solution de \textbf{gestion électronique du courrier} : tests d'acceptation, \textbf{validation des exigences métier}, accompagnement des utilisateurs finaux et \textbf{support au déploiement}
                      \end{itemize}
                    }
                    {Tests fonctionnels,UAT,Tests E2E,Recette,Mobile Money,GED}
  \emptySeparator
  \experience
    {Juin 2025}   {Ingénieur Logiciel Junior}{Société Générale Cameroun}{Douala}
    {Jan. 2025} {
                      \begin{itemize}
                        \item Conception et développement d'une \textbf{solution de gestion des anomalies opérationnelles} en milieu bancaire
                        \item \textbf{Réduction de 25\,\% du temps de traitement} des anomalies et \textbf{diminution de 35\,\% des erreurs manuelles}
                        \item Amélioration de la traçabilité des résolutions et suivi en temps réel des KPI opérationnels
                      \end{itemize}
                    }
                    {.NET Core,KPI,Milieu bancaire}
  \emptySeparator
  \experience
    {Fév. 2024}   {Développeur Logiciel}{Boissons du Cameroun}{Douala}
    {Sept. 2023} {
                      \begin{itemize}
                        \item Digitalisation du \textbf{processus d'analyse des causes profondes SCAT -- DNV} via une application web
                        \item \textbf{Réduction de 40\,\% du temps d'investigation} des incidents et amélioration de la qualité des rapports d'analyse
                      \end{itemize}
                    }
                    {Symfony,Flask,SCAT -- DNV}
  \emptySeparator
  \experience
    {Avr. 2023}   {Analyste Programmeur}{Perenco Rio del Rey}{Douala}
    {Jan. 2023} {
                      \begin{itemize}
                        \item Optimisation des \textbf{processus de contrôle qualité QHSE} par le développement d'une application mobile et de macros Excel
                        \item \textbf{Amélioration de 20\,\% de l'efficacité des audits} et meilleure conformité réglementaire
                      \end{itemize}
                    }
                    {Kotlin,VBA Excel,QHSE}
  \emptySeparator
  \experience
    {Oct. 2022}   {Développeur Back-End}{WinMoney Cameroun}{Yaoundé}
    {Juin 2022} {
                      \begin{itemize}
                        \item Implémentation d'une \textbf{API RESTful} pour une plateforme de création de contenus visuels
                        \item Support de \textbf{jusqu'à 500 requêtes/minute} avec un temps de réponse moyen de 300\,ms
                      \end{itemize}
                    }
                    {Spring Boot,API REST}
  \emptySeparator
  \experience
    {Oct. 2021}   {Développeur Back-End Junior}{Kiama S.A}{Yaoundé}
    {Juin 2021} {
                      \begin{itemize}
                        \item Mise en place d'un \textbf{système de recommandation multicanal (email/SMS)} pour \link{https://ekiosque.cm}{ekiosque.cm}
                        \item \textbf{Augmentation de 15\,\% de l'engagement utilisateur} et amélioration du taux de conversion
                      \end{itemize}
                    }
                    {Symfony,PHP-ML,Email/SMS}
\end{experiences}
		% Section expérience
	% Awesome Source CV LaTeX Template
%
% This template has been downloaded from:
% https://github.com/darwiin/awesome-neue-latex-cv
%
% Author:
% Christophe Roger
% Contenu adapté pour : Samuel SIKATI KENMOGNE
%
% Template license:
% CC BY-SA 4.0 (https://creativecommons.org/licenses/by-sa/4.0/)

%Section: Formation
\sectionTitle{Formation}{\faGraduationCap}

\begin{scholarship}
	\scholarshipentry{2022 -- 2025}
					{\textbf{Ingénieur Informatique, Génie Logiciel (BAC+5)} -- \textsc{Institut UCAC-ICAM} (Ex IST-AC), Douala\newline
					\footnotesize{Gestion de projet, conduite du changement, méthodes agiles (Scrum, Kanban), management des SI, entrepreneuriat. Architecture logicielle avancée, développement d'applications web et mobile, DevOps, modélisation (UML, SysML, BPM).}}
	\scholarshipentry{2019 -- 2022}
					{\textbf{Licence en Informatique, Génie Logiciel (BAC+3)} -- \textsc{Institut Africain d'Informatique}, Yaoundé\newline
					\footnotesize{Administration de bases de données (SQL, NoSQL), conception d'API, administration systèmes et réseaux (Windows/Linux).}}
\end{scholarship}

%Section: Certifications
\sectionTitle{Certifications}{\faCertificate}

\begin{scholarship}
	\scholarshipentry{11/2024}{\textbf{Python Essentials 2} -- Cisco Networking Academy}
	\scholarshipentry{10/2024}{\textbf{PCEP} -- Certified Entry-Level Python Programmer -- OpenEDG Python Institute}
	\scholarshipentry{04/2024}{\textbf{TOEIC} (800/990) -- ETS EMEA}
	\scholarshipentry{08/2023}{\textbf{Scrum Fundamentals Certified} -- VMEdu Inc}
\end{scholarship}
				% Section formation + certifications
	% Awesome Source CV LaTeX Template
%
% This template has been downloaded from:
% https://github.com/darwiin/awesome-neue-latex-cv
%
% Author:
% Christophe Roger
% Contenu adapté pour : Samuel SIKATI KENMOGNE
%
% Template license:
% CC BY-SA 4.0 (https://creativecommons.org/licenses/by-sa/4.0/)

%Section: Languages
\twocolumnsection
{\sectionTitle{Langues}{\faLanguage}
\begin{skills}
	\skill{Français}{5}
	\skill{Anglais}{4}
\end{skills}}
{\sectionTitle{Forces}{\faPlus}
\vspace{1em}
\begin{itemize}
	\item Adaptabilité
	\item Travail d'équipe
	\item Communication technique
	\item Résolution de problèmes
\end{itemize}
}
					% Section langues + forces
	% Awesome Source CV LaTeX Template
%
% This template has been downloaded from:
% https://github.com/darwiin/awesome-neue-latex-cv
%
% Author:
% Christophe Roger
% Contenu adapté pour : Samuel SIKATI KENMOGNE
%
% Template license:
% CC BY-SA 4.0 (https://creativecommons.org/licenses/by-sa/4.0/)

%Section: Centres d'intérêt
\sectionTitle{Centres d'intérêt}{\faHeart}

\begin{keywords}
	\keywordsentry{Technologie}{Veille technologique}
	\keywordsentry{Loisirs}{Football, Lecture, Jeux d'échecs}
\end{keywords}
				% Section intérêts
	%% Awesome Source CV LaTeX Template
%
% This template has been downloaded from:
% https://github.com/darwiin/awesome-neue-latex-cv
%
% Author:
% Christophe Roger
% Contenu adapté pour : Samuel SIKATI KENMOGNE
%
% Template license:
% CC BY-SA 4.0 (https://creativecommons.org/licenses/by-sa/4.0/)

%Section: Projets
% Aucun projet personnel n'est listé dans le CV actuel.
% Modèle prêt à l'emploi : décommenter % Awesome Source CV LaTeX Template
%
% This template has been downloaded from:
% https://github.com/darwiin/awesome-neue-latex-cv
%
% Author:
% Christophe Roger
% Contenu adapté pour : Samuel SIKATI KENMOGNE
%
% Template license:
% CC BY-SA 4.0 (https://creativecommons.org/licenses/by-sa/4.0/)

%Section: Projets
% Aucun projet personnel n'est listé dans le CV actuel.
% Modèle prêt à l'emploi : décommenter % Awesome Source CV LaTeX Template
%
% This template has been downloaded from:
% https://github.com/darwiin/awesome-neue-latex-cv
%
% Author:
% Christophe Roger
% Contenu adapté pour : Samuel SIKATI KENMOGNE
%
% Template license:
% CC BY-SA 4.0 (https://creativecommons.org/licenses/by-sa/4.0/)

%Section: Projets
% Aucun projet personnel n'est listé dans le CV actuel.
% Modèle prêt à l'emploi : décommenter \input{section_projets} dans cv.tex
% puis remplacer les champs entre < >.
\sectionTitle{Projets}{\faLaptop}

\begin{projects}
	\project
	{Site personnel}{<années>}
	{\website{https://samuelsikati.com}{samuelsikati.com} \github{sikatikenmogne}}
	{<Description du projet>}
	{<Techno 1>,<Techno 2>}
\end{projects}
 dans cv.tex
% puis remplacer les champs entre < >.
\sectionTitle{Projets}{\faLaptop}

\begin{projects}
	\project
	{Site personnel}{<années>}
	{\website{https://samuelsikati.com}{samuelsikati.com} \github{sikatikenmogne}}
	{<Description du projet>}
	{<Techno 1>,<Techno 2>}
\end{projects}
 dans cv.tex
% puis remplacer les champs entre < >.
\sectionTitle{Projets}{\faLaptop}

\begin{projects}
	\project
	{Site personnel}{<années>}
	{\website{https://samuelsikati.com}{samuelsikati.com} \github{sikatikenmogne}}
	{<Description du projet>}
	{<Techno 1>,<Techno 2>}
\end{projects}
				% Section projets (à compléter)
	%% Awesome Source CV LaTeX Template
%
% This template has been downloaded from:
% https://github.com/darwiin/awesome-neue-latex-cv
%
% Author:
% Christophe Roger
% Contenu adapté pour : Samuel SIKATI KENMOGNE
%
% Template license:
% CC BY-SA 4.0 (https://creativecommons.org/licenses/by-sa/4.0/)

%Section: Références
% Pour ajouter des référents, remplacer le paragraphe par l'environnement :
% \begin{referees}
% 	\referee{<Nom>}{<Poste>}{<Entreprise>}{<email>}{<téléphone>}
% \end{referees}
\sectionTitle{Références}{\faQuoteLeft}

\par{Disponibles sur demande.}
				% Section références
	%% Awesome Source CV LaTeX Template
%
% This template has been downloaded from:
% https://github.com/darwiin/awesome-neue-latex-cv
%
% Author:
% Lambtho12
%
% Reference section based on the template of LianTze Lim
% https://www.overleaf.com/latex/templates/a-customised-curve-cv/mvmbhkwsnmwv
%
% Template license:
% CC BY-SA 4.0 (https://creativecommons.org/licenses/by-sa/4.0/)

%Section: Publications
% Aucune publication dans le CV actuel : cette section est désactivée dans cv.tex.
% Pour l'activer :
%   1. renseigner vos publications dans my_publications.bib
%   2. décommenter \addbibresource{my_publications.bib} et % Awesome Source CV LaTeX Template
%
% This template has been downloaded from:
% https://github.com/darwiin/awesome-neue-latex-cv
%
% Author:
% Lambtho12
%
% Reference section based on the template of LianTze Lim
% https://www.overleaf.com/latex/templates/a-customised-curve-cv/mvmbhkwsnmwv
%
% Template license:
% CC BY-SA 4.0 (https://creativecommons.org/licenses/by-sa/4.0/)

%Section: Publications
% Aucune publication dans le CV actuel : cette section est désactivée dans cv.tex.
% Pour l'activer :
%   1. renseigner vos publications dans my_publications.bib
%   2. décommenter \addbibresource{my_publications.bib} et % Awesome Source CV LaTeX Template
%
% This template has been downloaded from:
% https://github.com/darwiin/awesome-neue-latex-cv
%
% Author:
% Lambtho12
%
% Reference section based on the template of LianTze Lim
% https://www.overleaf.com/latex/templates/a-customised-curve-cv/mvmbhkwsnmwv
%
% Template license:
% CC BY-SA 4.0 (https://creativecommons.org/licenses/by-sa/4.0/)

%Section: Publications
% Aucune publication dans le CV actuel : cette section est désactivée dans cv.tex.
% Pour l'activer :
%   1. renseigner vos publications dans my_publications.bib
%   2. décommenter \addbibresource{my_publications.bib} et \input{section_publications} dans cv.tex
\sectionTitle{Publications}{\faBook}

\begin{publications}
\nocite{*}

\printbibliography[heading=subbibliography,title={Articles},type=article]

\printbibliography[heading=subbibliography,title={Conférences},type=inproceedings]

\printbibliography[heading=subbibliography,title={Livres et chapitres},filter={booksandchapters}]

\end{publications}
 dans cv.tex
\sectionTitle{Publications}{\faBook}

\begin{publications}
\nocite{*}

\printbibliography[heading=subbibliography,title={Articles},type=article]

\printbibliography[heading=subbibliography,title={Conférences},type=inproceedings]

\printbibliography[heading=subbibliography,title={Livres et chapitres},filter={booksandchapters}]

\end{publications}
 dans cv.tex
\sectionTitle{Publications}{\faBook}

\begin{publications}
\nocite{*}

\printbibliography[heading=subbibliography,title={Articles},type=article]

\printbibliography[heading=subbibliography,title={Conférences},type=inproceedings]

\printbibliography[heading=subbibliography,title={Livres et chapitres},filter={booksandchapters}]

\end{publications}
			% Section publications (aucune pour le moment)

\end{document}
